\sect{Заключение}{conclusion}
%
В настоящей работе представлено единое математическое описание
моделей Hidden Markov Model и On-line Time Warping для задачи
отслеживания партитуры в реальном времени, а также архитектура
гибридного конвейера, объединяющего обе модели и снабжённого модулем
обучения с подкреплением для опережающего предсказания темпа на
горизонте $200$--$500$\,мс. Конвейер включает в себя обучаемый
свёрточный фронтенд, заменяющий гауссову функцию эмиссии HMM при
работе с аудиовходом, и систематический набор правил для обработки
граничных случаев~--- ферматы, орнаменты, пропуски, повторы и
октавные ошибки. Эксперименты на синтетическом датасете в стиле
прелюдий Шопена, на открытых корпусах MAPS и MSMD, а также на
собственном датасете \textit{CU-Concerto-2026}, собранном в
Центральном Университете, показали, что предложенная схема
снижает эффективную задержку аккомпанемента с $34$ до $19$\,мс и
улучшает F1-метрику с $0{,}947$ до $0{,}979$ по сравнению с базовой
HMM-конфигурацией.

Главный методологический вклад работы~--- демонстрация того, что
обучение с подкреплением полезно не как замена классического трекера,
а как отдельный опережающий слой, действующий поверх него. Такой
подход сохраняет вероятностную интерпретируемость скрытых марковских
моделей, преимущества OLTW по устойчивости к рассогласованию и
одновременно даёт инструмент адаптации темпа под индивидуальный
стиль исполнителя.

Дальнейшие исследования планируется развивать в трёх направлениях.
Первое~--- замена явного вознаграждения~\eqref{eq:rl-reward} на парные
предпочтения, собираемые от музыкантов консерватории, по
схеме~\cite{Christiano2017Preference}: это позволит включить в
обучение эстетические критерии, не сводимые к одной численной метрике.
Второе~--- формализация всего трекера как POMDP над
belief-состоянием HMM с агентом, выбирающим действия из дискретного
множества альтернатив~\cite{Kaelbling1998POMDP}. Третье~--- расширение
датасета \textit{CU-Concerto-2026} и репертуара поддерживаемых
произведений, а также пилотное внедрение в учебный процесс
Центрального Университета и партнёрских консерваторий.
